(1) 取$\psi_0(t)=A_0e^{iQA_0^2t}$，代入方程可见$\psi_t=iQA_0^2\psi_0$且$\psi_{xx}=0$，所以$i\psi_t+Q|\psi|^2\psi=-QA_0^2\psi_0+QA_0^2\psi_0=0$。该背景是理想无限均匀波，不代表有限能量局域孤子。
(2) 令$\psi=[A_0+u(x,t)+iv(x,t)]e^{iQA_0^2t}$，$u,v$为小实扰动。保留一阶得$u_t=-P v_{xx}$、$v_t=P u_{xx}+2QA_0^2u$。设$(u,v)\propto e^{i(Kx-\Omega t)}$，消元得到$\Omega^2=P^2K^4-2PQA_0^2K^2$，因此侧带是否出现虚频完全由色散与非线性符号组合决定。
(3) 若$PQ>0$，当$0<K^2<2QA_0^2/P$时$\Omega^2<0$，增长率$\sigma(K)=\sqrt{2PQA_0^2K^2-P^2K^4}$。对$K^2$求导得峰值位于$K_*^2=QA_0^2/P$，最大增长率$\sigma_{\max}=|Q|A_0^2$。如果$PQ<0$，所有$K$的$\Omega^2\ge0$，此类平面波线性调制稳定。
(4) 无损NLS有粒子数（光功率）守恒$\partial_t|\psi|^2+\partial_x J=0$，$J=2P\,\operatorname{Im}(\psi^*\psi_x)$，总积分不随时间改变。线性不稳定意味着侧带从背景重新分配强度，不意味着凭空增加总功率；当扰动达到$O(A_0)$，线性增长率失效，须用非线性演化或更完整的色散/损耗模型。
